\chapter{绪论}
\section{选题背景及意义}
随着大语言模型、前端框架 React、Vue 的兴起，
作为理论计算机历史悠久的分支之一编程语言理论\&设计实践（Programming Language Theory，PLT \& PL Design \& Implementation, PLDI）与相关领域类型论（Type Theory，TT）、范畴论（Category Theory）重新进入人们的视野。
作为其产物之一的函数式编程范式（functional programming）被活用于前端设计部分，
而另一产物定理证明器则与大模型结合，产生了 AI4MATH 等证明自动化项目。
尝试利用大模型解决奥林匹克数学问题的捷报频传，受到数学家陶哲轩等人的青睐。
然而因其理论程度高（冷门）、对学生纯数学水平要求高的特点，许多大学并未开设该课程，
推广 PLT、PLDI、类型论和范畴论的教育与学界所倡导的编程实践仍然是一项挑战。

本文旨在探索一种同时面向工程化程序开发与形式化验证的语言与编译工具链原型，
同时获得可运行的编译器产出与可用的证明重写器，作为本科生或研究生科研训练与工程原型，
从而成为教学、实验与工程验证的良好平台。
这一原型分为两个子系统（Tigris 语言与 Tigrisd 定理证明器，两者合称 TinyTheorem 解决方案）按序实现：

在程序语言实践层面，本文继承并扩展 Hindley–Milner（HM）类型系统的工程优势（类型推断、简单语义、可扩展的类型类与谓词），通过在工程化实践与更强语言表达力之间进行取舍，创新性地设计出高阶多态、蕴含的谓词约束、严格控制的高阶类型 HKT 构造等语言（类型系统）特性，系统化提升可读、可维护的中型编译器构建能力。产物包括 Tigris 语言、其形式化语言规格（language specifications）和一个可输出多层次结果的编译器、一个基于 IR 的解释器与 REPL。这一子系统使用 Lean 实现；

在形式化验证侧，本文设计了简化的依赖类型（dependent types）内核，
以精简版的依赖类型\lambdacalc（$\lambda$-Calculus）为计算模型，并采用双向（bidirectional）类型检查，
在单一宇宙层级（1-level type theory）下实现蕴含的（impredicative）特性以降低内核复杂度\upcite{norell2007towards,loh2010tutorial}，
用于支持基本的证明对象与重写推理，
产物包括 Tigrisd 语言（定理证明器）、其形式化语言规格和一个可用于定理证明的命令行前端。这一子系统使用 Haskell 实现。

本文借助“解耦”思路，以分离式架构将两个子系统并行推进：编程子系统坚持HM的相对弱的表达力与可编译性；证明子系统把依赖类型的强大表达力局限在更小而可控的内核上，从而避免将全依赖类型的复杂度与代价强加于普通程序员的日常开发环节。
而项目本身，在不牺牲工程实践友好性的前提下，提供一体化但解耦的“语言 + 证明”原型，既便于教学与课程项目落地：如类型推断、闭包转换与续元传递（continuation-passing）化、函数式编程与依赖类型编程实践，
也为进一步研究（如更强的多态、约束求解、编译优化与证明自动化）提供坚实基础与统一框架。

\section{研究现状与文献综述}
以下叙述本项目涉及的研究的现状。为了体现本项目的特殊之处，在每一部分的最后一段，若本文的实现与之有创新，我们加以简要说明。

需要明确的先验认识是：定理证明器也是编程语言，且一般为纯函数式语言，
因为函数式语言在数理逻辑/数学上的合理性（soundness）最好，其中又以纯函数式为甚。
举例来说，数学的函数都没有副作用，这一点在 Haskell 语言中得到体现，除非是
单子上（Monadic，是受范畴论同名代数结构启发，函数式编程中用于建模副作用概念）的函数\footnote{
  而即便是这样的函数，
  除去输入输出、多线程、STM（Software transactional memory）这些真正与副作用相关的 \texttt{IO}、\texttt{ST} Monad 之外，其余都可以用纯函数\emph{模拟}出来，本文的实现亦完全符合这一设计模式，因为 Haskell 和 Lean 都是
  纯函数式语言。
}。
\subsection{Hindley-Milner 类型系统} 由 Hindley、Milner 等奠基并由 Damas–Milner 系统化，形成“具有主类型/类型模式（principal type-schemes）”与“算法W（Algorithm W）”的经典框架\upcite{hindley1969principal,milner1978theory,damas1982principal}。其核心优势为：对 $\lambda$-项的完全/全局类型推断（complete/global type inference）、良好的语义基础与实现简洁性。该系统的后续工作围绕“向多态更强方向扩展”与“保持推断可行性”展开：
\paragraph{类型类（typeclass）与限定类型（qualified types）} 由 Wadler、Blott、Jones 等提出与系统化，提供了对重载与约束多态的统一抽象，使得字典传递成为行之有效的编译期（compile-time）繁饰（elaboration）设计策略\upcite{wadler1989make,jones1993system,jones1995functional}。
\paragraph{高阶多态（higher-rank polymorphism）的推断} 由 Peyton Jones 等提出，Dunfield–Krishnaswami 等人则提供了“完备且易实现的双向类型检查”解决方案，在实践中往往以“显式标注 + 局部推断 + 刚性检查”的折中方式落地\upcite{jones2007practical,dunfield2013complete}。
\paragraph{局部类型推断（local type inference）与双向类型检查} 由 Pierce 提出并发展，则取强表达力与用户心智负担的中间地带：
仅在需要时（例如函数声明，或2-阶及以上时）通过少量标注引导推断，兼顾可读性与可推断性\upcite{dunfield2013complete,pierce2000local}。

\paragraph{本文的创新点} 在本文的“编程子系统”中，我们借鉴上述成果，采用“扩展HM + 谓词（predicates） + 类型类 + 严格一阶类型（strict first-class） + 头等泛型（first-class polymorphism）”的设计，并配以“内嵌一阶模式（Scheme）方案 + Skolem 刚性检查 + 承担字典投影（dictionary projection）的System F繁饰”的实现策略，使其既与 Haskell 传统保持一致性，又便于在原型编译器框架下实现与调试。

\subsection{语法与解析（parsing）} 函数式语言（仅指 Milner 等人开创的 ML 系语言）通常具有相仿
的语法结构。其中著名者包括 Haskell、OCaml、Standard ML；在定理证明器部分则又包括 Coq （现称 Rocq）、Lean、Agda、Isabelle/HOL 与 F*. 函数式语言语法的特点是灵活与统一，
既允许多种多样的用户自定义运算符（例如使用 LR(1) ocamlyacc 实现的 OCaml），更甚者
还支持强大的元编程（metaprogramming）功能，例如 GHC 的 TemplateHaskell，OCaml 的 Pre-Processor Extension （PPX）和 Lean 的 \texttt{syntax}、\texttt{macro}、\texttt{macro\_rules} 等语法构造：
这一特性在定理证明器中尤为明显，因其对数学证明的需求；统一性则体现在函数式语言对语法糖（syntactic sugar）的缺失，得益于这些语言
的面向表达式的特性（expression-oriented），其表层语法（surface syntax）具有强通用性，
且统一使用含类型（typed）\lambdacalc（或其变体，例如定理证明器则使用依赖类型的版本）作为它们的核心代数（core calculus）。从这个角度，也可以说函数式语言\emph{几乎只有}语法糖，因为其表层语法全部被解糖（desugar）到前述的核心代数中，但为了和工业语言特殊对待的“语法糖”区分，我们不称其为语法糖。

\paragraph{本文的创新点} 这些特性，除元编程之外，都由 Tigris 继承。我们设计的表层语法综合了 Lean、OCaml、Haskell 的长处，
并通过受范畴论启发的单子式\footnote{属于函数式编程的设计模式的一部分。}
的语法分析组合子（monadic parser combinator）\upcite{frost1989constructing,Hutton_1992,frost2008parser}
保证解析器的模块化和可维护性，同时利用 Pratt 解析技术实现\upcite{pratt1973top}“用户可定义运算符 + 优先级/结合性”的高可扩展解析能力。
Pratt 自顶向下运算符优先级解析法以其“实现简洁、扩展性强”而被广泛采用，
适合教学与原型构建。结合“分派/分组式 let 语法糖”、
模式匹配（pattern matching）与记录（record）/归纳（代数）数据类型（inductive types/algebraic datatypes）等语言构件，
可覆盖函数式编程主流设计模式。

\subsection{类型类的繁饰与实例解析（instance resolution）} 基于 Wadler-Blott 的类型类思想，
将临时（ad-hoc）多态（例如 C、Java 的函数重载）以“字典”抽象归结为参数化（parametric）多态的系统化繁饰；
Peyton Jones 的合格类型与“构造子类”进一步将类型构造子（如 Functor、Applicative、Monad 等经典纯函数式设计模式\upcite{moggi1991notions}）纳入框架。
在实际实现中，Haskell 主流实现之一的 GHC 的 OutsideIn(X)\upcite{vytiniotis2011outsidein} 体系及其后续演化提出了以约束为中心的求解与实例选择流程，
并辅以 Skolem 刚性变量（rigids）、泛化（generalization）/实例化策略来控制推断过程的完备性与可预测性。

\paragraph{本文的创新点} 本文采用与之相容的思路，但简化了 OutsideIn(X) 体系：
System F 的繁饰以字典参数（项层）+ 类型抽象（类型项层）联合刻画；
刚性检查通过在注解与应用位点进行 Skolem 化检查，防止不当的跨层次统一；
实例解析采用“基于头部匹配 + 递归解析上下文谓词”的策略（参照 OutsideIn 风格），并结合模式匹配决策树优化编译\upcite{maranget2008compiling}。

\subsection{IR 设计与编译后端} 在编译管线方面，
研究社区已就“System F 繁饰 $\to$ Administrative-Normal Form（ANF）$\to$ 闭包转换（closure conversion）$\to$ CPS 化”
的路径形成较为成熟的共识与方法论。这其中，
Girard提出的System F作为多态\lambdacalc 的标准中间表示，适合承载类型抽象与字典化后的多态程序\upcite{reynolds1983types}；
Flanagan、Appel 等人在 ANF 与 CPS 化的经典工作明确了将控制流与求值顺序外显（explicit）化的工程路径\upcite{appel2007compiling,flanagan1993essence}，
利于优化与目标后端对接。

\paragraph{本文的创新点}我们在此基础上，进一步设计了“Lambda IR（ANF 变体）$\to$ CPS IR”的两段式IR，
并最终落地到 Common Lisp\upcite{steele1990common}以获得成熟运行时、交互环境、外部函数接口（foreign function interface，FFI）的支持。

\subsection{依赖类型内核与双向检查} 在证明子系统方面，
Lean、Coq、Agda\upcite{moura2021lean,barras1999coq,bove2009brief}分别代表了现代依赖类型定理证明器的不同取向：
Lean 与 Coq 的理论基础一致（Calculus of Inductive Constructions, CIC）\upcite{coquand1986calculus,Dybjer_1991,10.1007/BF01211308}，但前者强调工程机制与元编程生态，
后者强调严谨公理化\upcite{10.1007/BFb0037116,10.1007/BFb0040259}\footnote{
  前者与后者相比，提供了一些现代逻辑认为不可靠的公理，例如集合论的选择公理（Axiom of Choice），以及由之导出的排中律（Law of Excluded Middle），因而受人诟病；
  而考虑到编程需要，其在停机检查 termination checking 时也较松，允许 partial 和 unsafe. 因而本项目选择用其作为编程语言。
}，
而 Agda 则强调程序化证明与可视化交互，
且主要用于建模（modeling）类型论研究前沿的同伦类型论（Homotopy Type Theory）与立方类型论（Cubical Type Theory）。
而它们的双向检查部分，基本来源于前述编程语言部分类型系统的实现。

\paragraph{本文的创新点}本项目采用“简化依赖类型 + 单一宇宙 + 双向检查”的取舍，参考教学与研究社区中的“pi-forall”、“LambdaPi”等项目经验，
避免把复杂的宇宙层级、多层归纳族与高阶等式推理一次性纳入，实现“可运行、可实验”的证明原型\upcite{loh2010tutorial, brady2005practical}。通过与编程子系统解耦，
用户在不涉及证明时可完全在 HM 子系统内获得流畅的开发体验，
在需要轻量证明与重写时又可切换到依赖类型内核。

\section{主要研究内容、预期目标与创新点}
\subsection{研究的主要内容与设计思路}

在语言前端与解析上，我们将实现 Lean/OCaml/Haskell 风格交融的函数式语法，采用 Pratt 解析构建可扩展的表达式优先级框架，支持用户自定义操作符（优先级/结合性）与冲突检查。表层语法大体依照 OCaml 风格，但参考 Lean 的习惯写法提供语法糖和 Haskell 的写法提供部分语法结构（language construct）与模式匹配、记录/ADT、头等抽象、类型标注（type annotation/ascription）等语法。

在类型系统实现上，构建“扩展 HM + 谓词 + 类型类 + 第一阶 HKT”类型系统；支持 rank-1 方案（TSch）内嵌与字典化繁饰；在语法树（abstract syntax tree, AST）的 \texttt{Ascribe/Var/App} 等关键位点使用 Skolem 刚性检查与rank-1 方案实例（instantiation）；提供局部双向检查与必要的显式标注通道。在类型求解上：采用一阶合一（unification）的方式求解类型、参照“字典变量在作用域内匹配/查找”的 OutsideIn 式思路，维护“已访问实例”集合避免出现环导致系统崩溃（diverge）；在类型制导（type-directed）语法上，我们令类型检查器（typechecker）在需要时（譬如实例解析、解糖时的易歧义处）要求显式类型注解，提升错误定位与歧义时的可预期性。
在IR降级上，我们繁饰表层语法到 System F：新增 Proj（字典投影）等节点便于优化与降级；将类型抽象与字典参数显式化，提供\texttt{TyLam/TyApp/Proj}等 System F IR 结点，并配套$\beta/\eta$-化、减少生成的 IR 噪声并向后续步骤暴露更多的优化机会，同时保证后续变换的可维护性；

在后端 IR 与编译上，我们的管线是 Lambda IR $\to$ Closure Conversion $\to$ CPS Transformation $\to$ Common Lisp；在IR层实现常量折叠（constant folding）、拷贝传播（copy propagation）、无用代码消除（dead code elimination）、尾调用（tailcall/tail-recursive）化、投影/选择消去的源优化，最后在CPS层显式表达控制转移并直接下降到目标 Common Lisp 代码。
模式匹配检查与编译上：在类型全数推断出来后，我们参照 Maranget 的模式匹配冗余/欠缺检查算法，针对缺少的/多余的模式分支向用户发出警告，并针对具体情况向用户提出一种缺失的分支/立即删除多余分支；在System F 繁饰之后，我们根据其基于决策树的编译算法，生成良好分支。

在证明内核部分，我们采用单一宇宙、蕴含式的依赖类型\lambdacalc，并采用双向类型检查，降低判定复杂度与术语成本；
参考标准的等式重写与项代换理论（Baader\&Nipkow）\upcite{baader1998term}提供项重写式（term-rewriting）解释器以支持基本等式推理与程序化证明脚本；利用 Haskell 在学术上的丰富生态，通过 Unbound 等库实现对于\lambdaterm 的操作，变换，以及变量约束的判断；
引入 Peyton Jones 提出的 Weak Head 范式（Weak Head Normal Form, WHNF）\upcite{peyton1987implementation}用于控制消去（reduction），将其与定义性等号（definitional equality）的实现结合，避免在判断两个项是否相等时的无用计算、消去、展开（unfolding），提升定理证明性能。

最后，在实验与评测上，我们以 Functor/Applicative/Monad 等典型代数结构暨函数式编程设计模式为样例，验证实例解析、字典化与投影优化的正确性；以典型函数式程序映射（map）、折叠（fold）考察编译后端的性能与优化效果；
以基本逻辑定理、数学结构的编码作为证明例子，验证证明内核的可用性。

课题整体采用“自顶向下的最小可用（MVP）+ 逐步精炼”的策略。首先保证语言前端、类型系统核心路径、编译与运行链路闭环（从 Tigris 源到 Common Lisp 目标可运行），再迭代增强类型系统细节、实例解析策略与 IR 优化规则。编程与证明两子系统以统一仓库并行推进，但保持边界清晰、依赖方向单纯：编程侧产物可独立使用，
证明侧可以通过两系统相似的表层语法无痛迁移。

\subsection{研究的预期目标}
本课题包含三个目标（要注意，这是最低限度的目标，实际上我们实现的功能并不止与此，见正文）：
{\setlength{\parskip}{-1em}
  \paragraph{其一} 一个可运行的原型编译器与解释器：可编译到 Common Lisp（SBCL）或直接在 REPL 交互下运行；可对带类型类与 HKT 的示例进行正确繁饰求值。
  \paragraph{其二} 完整的IR管线与可视化输出（基于 Wadler 的 Pretty Printing 算法）：便于教学与调试；可展示从System F\upcite{girard1986system,harper2016practical} 到 Lambda IR、CPS IR 的逐步变换。
  \paragraph{其三} 证明侧最小可用内核：支持基本\lambdaterm、依值积/和（dependent product/coproduct, $\Pi/\Sigma$）与等式重写的简单证明脚本；在实际例子（如代数结构的基本定理）中可用。
}
\section{论述结构安排}\label{arrange}
本文主要对课题 TinyTheorem 的两个子系统（Tigris、Tigrisd）做出描述和形式化说明，以规范语言行为；
同时论述语言设计与实现上的考量、优势，以及回首复查时发现的设计失误并向读者介绍两系统的使用方式。
此外，本文不仅提倡函数式编程和依赖类型编程范式的使用，同时也是其实践者。
我们还将提供一个运用于本项目的依赖类型编程的实例作为范例，并向读者介绍其思考方式、安全性与优点。
具体来说，除去绪论，本文的正文部分可以分为六大部分：

\paragraph{第一部分}以语言明细/报告的（参照 OCaml Manual 或 Haskell Language Report）\upcite{marlow2010haskell,leroy2024ocaml}口吻用自然语言
向读者全面叙述 Tigris 语言的设计以及其各部分行为，
同时亦作为读者了解 Tigris 如何使用、如何工作的主要途径。因为编译是线性过程，自然，
我们跳出适合一般、多层次软件的描述框架，
按最直观且最符合逻辑的编程语言报告的格式，根据\emph{语言对象}的原子程度和编译
器各 pass 的顺序作为叙述顺序，具体是：

\begin{center}
  \emph{词法惯例、值、辨识符、类型表达式、（模式匹配的）模式、表达式/定义语句、类型定义}
\end{center}
在本部分，代码实现上的论述主要包含词法、句法分析，同时简单涉及类型检查器，即编译器的前端。这是因为类型系统和后续的编译部分和后端实现起来较为复杂，所以由它们将在后文各自独立的部分进行叙述。
同时，我们认为学术语言，特别是函数式语言和工业语言，特别是常见的面向对象语言在设计、使用、实现上都有诸多不同
（例如函数式语言通常是值导向的，并且具有统一的计算模型，这在工业语言中仍然较少见）。
为了避免犯\emph{全知视角}的写作错误，我们仍然从最基本的词法惯例开始叙述。
此外，这一部分亦作为读者了解函数式编程语言的窗口，因为 Tigris 语言是纯函数式的。

\paragraph{第二部分}专门用于给出正式的 Tigris 语言的类型系统的形式化表述。
这是因为类型系统、类型论在 PL/CAV/纯数学学界，是许多重要理论的基础，十分受到重视。我们仿照 PL 论文惯例，
利用 OTT\upcite{sewell2010ott} 给出严谨的 Tigris 的核心代数及其定义与判断规则（judgemental rules）、推断规则（inference rules）。
这些规则指导、规范着 Tigris 类型系统的实现，方便读者检查类型系统设计的 soundness，但阅读它们则不像阅读代码一样枯燥无味，但需要读者接受过充分的 PL/类型论的训练。
此外，
这一部分将更加深入的探讨 Tigris 类型系统、类型类与实例解析算法的实现及性质，最后将叙述 System F 繁饰的细节及
与之相关的一些编译期优化。

\paragraph{第三部分}对 Tigris 编译器的中后部分的实现与设计进行叙述，
包括 IR 设计、优化与生成、闭包转换、CPS 转换、后端代码生成、运行时系统及 FFI。
同时，还将涉及周边工具，如编译器命令行前端的设计使用，以及基于 IR 的解析器/REPL 的实现。
此外，在描述 CPS 转换时，我们给出正式的指称语义（denotational semantics），这些数学公式规定了 CPS 转换的对象
与结果，但需要读者受充分的 PL 训练。
这一部分应当是外行人在接触编译原理时首先想到的，其重要性不言而喻，且体现了\emph{编译}的字面义：
没有了这些管线，使用 Tigris 编写的程序终究无法工作（或者说，只能通过最直接的求值器工作）。

\paragraph{第四部分}撷取 Tigris 实现过程中出现的一个依赖类型编程范例，
以身作则向读者倡导、帮助读者理解这一编程范式的思考过程和优势。
可以看作是对依赖类型编程的一个简单描述，或一份教程。同时这一部分还起到承上启下的作用：
为读者提供即将叙述的 Tigrisd 一些先验知识，使读者能够更确切地理解后者的一些设计。

\paragraph{第五部分}主要给出 Tigrisd 定理证明器的自然语言介绍以及正式的形式化描述与实现细节，配套一个简单实例演示如何使用。由于 Tigrisd 在表层语法上与 Tigris 相仿，而且采用相似的
parser combinator 的句法分析实现方式，我们主要对两者的不同之处、与定理证明器特有的部分进行描述，主要包括
索引族（indexed family）\upcite{MARTINLOF1971179,10.1007/3-540-52335-9_47,10.1007/BF01211308}定义、函数、模式匹配与依赖积的消去（elimination）、证明/类型无关性、
内置命令以及类型宇宙（universe）。而最重要的 Tigrisd 内核
（kernel, 包含 typechecker 和类型论，依惯例称定理证明器的类型系统为类型论）、
与实现相关的 WHNF 和定义性相等则使用与第二部分非常相近的格式通过 OTT 给出形式化描述。
对内核的形式化描述规定了定理证明器的行为与计算法则（computational rule）。
阅读这一部分需要读者接受过充分的PL/类型论的训练。

\paragraph{第六部分}总结本课题并作展望，指出本项目以后可能的研究和发展方向，并提出一些改进与语言功能、配套编辑器插件上的增加。

\paragraph{读者须知} 最后，读者需要注意到，由于本课题自身属于编程语言设计与实现，正文中不可避免的会出现 Tigris 和 Tigrisd 的代码。
我们希望读者不要默认其枯燥无味，因为这些代码是\emph{样例}，而并非用于凑字数的具体实现的源代码；
它们用于形象地演示当前正在描述的问题。对于任何接触过编程的读者来说，我们认为相比于用自然语言描述问题，
使用样例代码可以更清晰的展现问题。同时，为了方便读者分辨样例代码与实现的源代码，在排版上\underline{对前者我们加浅蓝色边框}。

本文由英文原文翻译而来，表达上可能有欧化之处；文档由 \XeLaTeX 排版，参照数科院模板作了修改。

\section{本章小结}
本章从研究现状出发，介绍了 \Tigris、\Tigrisd 的理论基础，
并将学界的研究成果同我们的实现做比较，提出创新。
此外，本章的“论述结构安排”一节（\ref{arrange}）对每一章都做了全面、简洁的规划与介绍，读者应当重视。

TinyTheorem 解决方案的思路总的来说是避繁就简的，在学术功能与工业惯例间找寻平衡点，
以作为教育工具而不使学生感到无所适从。
